\begin{abstract}
Este trabajo aborda el problema de clasificación automática de imágenes de flores utilizando técnicas de aprendizaje automático y redes neuronales profundas. Se utiliza el conjunto de datos de la competencia \textit{Petals to the Metal -- Flower Classification on TPU} de Kaggle, compuesto por 23,847 imágenes RGB distribuidas entre 104 clases distintas de flores. El conjunto presenta variaciones importantes en tamaño, color, iluminación, orientación y composición de las imágenes, además de un desbalance en la cantidad de ejemplos disponibles para cada clase.

La tarea consiste en desarrollar un modelo capaz de identificar correctamente la clase de una flor a partir de una imagen. Debido a la similitud visual entre algunas categorías y a la variabilidad dentro de una misma clase, se considera principalmente el uso de redes neuronales convolucionales (CNN), ya que permiten extraer automáticamente características relevantes como formas, bordes, colores y patrones visuales. También se consideran arquitecturas previamente desarrolladas, como ResNet, AlexNet e Inception, así como métodos tradicionales de aprendizaje automático, entre ellos KNN, bosques aleatorios y máquinas de vectores de soporte, con el objetivo de comparar distintas alternativas para la solución del problema.

Para evaluar el desempeño se considera especialmente la métrica F1 macro utilizada por Kaggle, ya que asigna la misma importancia a cada una de las 104 clases independientemente de la cantidad de imágenes disponibles. De esta forma, el proyecto busca seleccionar y posteriormente implementar una estrategia de clasificación que logre un buen desempeño general sin favorecer únicamente a las clases con mayor cantidad de ejemplos.
\end{abstract}

\begin{keywords}
Clasificación de imágenes \sep
Redes neuronales convolucionales \sep
Aprendizaje profundo \sep
F1 macro \sep
Kaggle
\end{keywords}
